% Generated draft: requires mathematical and source review.
\documentclass[11pt]{article}
\usepackage[margin=25mm]{geometry}
\usepackage{fontspec}
\setmainfont{Latin Modern Roman}
\usepackage{amsmath,amssymb,mathtools,unicode-math}
\setmathfont{Latin Modern Math}
\usepackage{enumitem,needspace,xurl,hyperref}
\setlength{\parindent}{0pt}
\setlength{\parskip}{5pt}
\setlength{\emergencystretch}{3em}
\newcommand{\N}{\mathbb{N}}
\newcommand{\Z}{\mathbb{Z}}
\newcommand{\Q}{\mathbb{Q}}
\newcommand{\R}{\mathbb{R}}
\newcommand{\C}{\mathbb{C}}
\newcommand{\F}{\mathbb{F}}
\newcommand{\A}{\mathbb{A}}
\newcommand{\PP}{\mathbb{P}}
\newcommand{\E}{\mathbb{E}}
\newcommand{\eps}{\varepsilon}
\newcommand{\dd}{\,\mathrm{d}}
\begin{document}
\section*{UnsolvedMath ID: 701011}
\textbf{Source problem code:} catalogue entry
\subsection{Definitions and mathematical statement}
Work in the language of groups. A quasi-identity is a universally quantified implication
\[
\forall x_1,\ldots,x_r\quad
\left(\bigwedge_{i=1}^{m}u_i(x_1,\ldots,x_r)=v_i(x_1,\ldots,x_r)\right)
\Longrightarrow u(x_1,\ldots,x_r)=v(x_1,\ldots,x_r),
\]
where all displayed terms are group terms and the conjunction may be empty. A quasivariety of groups is a class axiomatized by quasi-identities.

Let \(\mathcal N\) be the class of all torsion-free nilpotent groups, without a fixed bound on nilpotency class. Here torsion-free means that \(g^n=1\), for \(n\geq 1\), implies \(g=1\). Nilpotence means that \(\gamma_{c+1}(G)=\{1\}\) for some finite integer \(c\geq 1\), where \(\gamma_1(G)=G\) and \(\gamma_{j+1}(G)=[\gamma_j(G),G]\). Let
\[
\mathcal Q=\operatorname{Qvar}(\mathcal N)
\]
be the smallest quasivariety containing \(\mathcal N\); equivalently, \(\mathcal Q\) consists of the groups satisfying every quasi-identity valid in every member of \(\mathcal N\).

Interpret the source's term ``direct wreath product'' in its usual unrestricted sense. For a group \(G\), let \(G^{\Z}=\prod_{n\in\Z}G\) be the group of all functions \(\Z\to G\), with pointwise multiplication and no finite-support restriction. For \(k\in\Z\), define \(\tau_k(f)(n)=f(n-k)\). The direct wreath product with the infinite cyclic group is
\[
G\wr\Z=G^{\Z}\rtimes_{\tau}\Z,
\qquad
(f,k)(h,\ell)=\bigl(f\,\tau_k(h),k+\ell\bigr).
\]
A quasivariety \(\mathcal M\) is closed under direct \(\Z\)-wreath products if \(G\wr\Z\in\mathcal M\) for every \(G\in\mathcal M\).

The question proposed by A. I. Budkin is whether
\[
\forall G\in\mathcal Q,\qquad G\wr\Z\in\mathcal Q.
\]
The supplied registry labels the problem open. Its accompanying discussion reports an explicitly unverified AI claim of an affirmative solution; this draft does not treat that claim as an established resolution.
\subsection{Short English statement}
Consider the quasivariety generated by all torsion-free nilpotent groups, allowing every finite nilpotency class. Must the unrestricted wreath product \(G\wr\Z\) belong to this quasivariety whenever \(G\) does? The supplied open label and the reported unverified affirmative AI claim are retained as source qualifications.
\subsection{Sources}
\begin{itemize}
\item Local UnsolvedMath registry, supplied complete record with permanent ID \(701011\), exact source problem code \texttt{catalogue entry}, title ``Kourovka Notebook Problem 14.26.'' The record attributes the problem to A. I. Budkin, proposed in 1999.
\item The supplied source attribution is \emph{The Kourovka Notebook (Unsolved Problems in Group Theory)}, 14th issue (1999), Problem 14.26, as included in the 21st issue edition (2026), \texttt{arXiv:1401.0300v48}. The supplied collection description names E. I. Khukhro and V. D. Mazurov as editors of Issues 1--21 (1965--2026).
\item The supplied discussion reports an explicitly ``unverified AI solution'' claiming an affirmative answer, attributed to M. Kinyon and J. Urban, \emph{Preprint}, 2026, with the supplied link \url{https://kourovkanotebookorg.wordpress.com/wp-content/uploads/2026/09/kourovka-mathematics.pdf}. No separate preprint title is supplied.
\item The supplied status review is dated \texttt{2026-10-16}, retains the canonical status ``open,'' and classifies the evidence as \texttt{unrefereed\_claim}. It says that the October 2026 21st issue, \texttt{arXiv:1401.0300v48}, marks the problem with \(\#\) for the unverified AI solution claim. This is supplied dated curation, not an independently verified literature assessment.
\end{itemize}
\end{document}
